% Propietario conservado: /Users/ruben/Documents/ChatGPT/jueces y controles/REVISION_ARGUMENTAL_APERTURAS_CIERRES_20260915/delivery/ARTICLE_V_ES_ARGUMENTAL_20260915/manuscrito/sections/10_realizacion_genealogica_evolucion.tex
% Recuperación documental para el artículo V; RESULTADO_RECUPERADO.
% Propietario: /Users/ruben/Documents/New project/output/REVISION_EDITORIAL_HMT_MD_20260905/01_FUENTES/INTEGRAL/tree/New project/output/ENSAMBLADO_CAUSAL_RESTAURADO_HMT_MD_20260905/fuente/manuscrito/integracion_83/parte_iv/body/B011_ch53_trayectorias_cuanticas_6b08e0e532c9_04ab_postulados_cuanticos_genealogicos_rev2.tex
% Líneas de origen: 4--95, 178--206.
% Sólo se adaptan las etiquetas de referencia y, cuando procede, el nivel de títulos.
% La matriz de recuperación conserva la correspondencia y las dependencias.
\section{Reconstrucción de los principios cuánticos desde el estado genealógico}
\label{v:v:rec:sec:postulados-cuanticos-genealogicos-rev2}
\index{mecánica cuántica!reconstrucción genealógica}
\index{estado!realización de Hilbert}

La formulación de Hilbert organiza la mecánica cuántica mediante estados,
observables autoadjuntos, evolución unitaria, composición tensorial y una
regla probabilística para los resultados~\cite{v:vonNeumann1932}. HMT--MD
construye una capa anterior a esa representación: una sección compatible,
su ruta observable, la memoria del transporte y el lector que publica un
resultado. El espacio de Hilbert no se reemplaza; aparece como codominio de
una realización lineal declarada.

La secuencia tipada es
\[
 \boxed{\begin{aligned}
 \text{sección genealógica}
 &\longrightarrow \text{realización lineal}
 \longrightarrow \text{fase y superposición}\\
 &\longrightarrow \text{observable y evolución}
 \longrightarrow \text{probabilidad y medición}
 \end{aligned}}
\]
Cada flecha conserva su dominio. En particular, el vector de estado no
sustituye a la sección TPK, del mismo modo que el valor real no sustituye a
la genealogía numérica que lo produce.

\subsection{Estado genealógico y realización de Hilbert}
\label{v:v:rec:realizacion_evolucion:section:02}

\begin{definition}[Estado genealógico realizado]
\label{v:v:rec:realizacion_evolucion:statement:01}
Un estado genealógico es una quíntupla
\begin{equation}
 \mathfrak G=(x_\infty,\Gamma,\mathcal F,\mathcal M,\mathcal R),
 \label{v:v:rec:eq:estado-genealogico-realizado-rev2}
\end{equation}
donde \(x_\infty\in\Omega_\infty\) es una sección compatible,
\(\Gamma\) su ruta observable, \(\mathcal F\) la familia de realizaciones,
\(\mathcal M\) la memoria de los levantamientos y \(\mathcal R\) el lector
declarado. Sea \(\mathsf P_{\rm TPK}\) la categoría cuyos objetos son las
secciones TPK preparadas, finitas o coinductivas, y cuyos morfismos son las
rutas registradas compatibles con sus datos de frontera. En particular, la
restricción de profundidad es un morfismo
\[
 r_{n+1,n}:x_{n+1}\longrightarrow x_n.
\]
Sea \(\mathsf{Hilb}_{\mathbb R,J}\) la categoría cuyos objetos
son pares \((\mathcal H_{\mathbb R},J)\), con \(\mathcal H_{\mathbb R}\)
espacio de Hilbert real y \(J\) ortogonal, \(J^2=-I\), y cuyos morfismos son
aplicaciones lineales reales acotadas \(T\) que entrelazan las estructuras
complejas, \(TJ_1=J_2T\). Una realización cuántica es un funtor
\[
 \mathcal Q:\mathsf P_{\rm TPK}\longrightarrow
 \mathsf{Hilb}_{\mathbb R,J},
\]
y una \emph{realización apuntada} es el par \((\mathcal Q,\Psi)\) que asigna
a cada sección preparada \(x\) un vector no nulo
\(\Psi_x\in\mathcal Q(x)\), compatible con las restricciones del sistema
inverso.
\end{definition}

Con la convención lineal en el segundo argumento, la forma hermítica asociada
se escribe
\[
 \langle u,v\rangle_{\mathbb C}
 =\langle u,v\rangle_{\mathbb R}
  -i\langle u,J_E v\rangle_{\mathbb R}.
\]
Elegir la carta \(J_E\leftrightarrow i\) realiza la estructura compleja que
en el núcleo formal y su desarrollo se construye como raíz interna de \(-I\). Para prefijos de
profundidad \(n\), pongamos
\(\mathcal H_n:=\mathcal Q(x_n)\) y
\[
 p_{n+1,n}:=\mathcal Q(r_{n+1,n}):
 \mathcal H_{n+1}\longrightarrow\mathcal H_n.
\]
La compatibilidad entre restricción y realización exige
\begin{equation}
 p_{n+1,n}\Psi_{n+1}(x_{n+1})=\Psi_n(x_n).
 \label{v:v:rec:eq:compatibilidad-realizacion-hilbert-rev2}
\end{equation}
Si \(r_{\infty,n}:x_\infty\to x_n\) es la restricción coinductiva y
\(p_n:=\mathcal Q(r_{\infty,n}):\mathcal H_\infty\to\mathcal H_n\), un
límite \(\mathcal H_\infty\) realiza la sección completa cuando
\(p_n\Psi_\infty(x_\infty)=\Psi_n(x_n)\) en cada profundidad.

\begin{principle}[Estado cuántico realizado]
\label{v:v:rec:realizacion_evolucion:statement:02}
El estado puro asociado a una sección se representa por el rayo del vector
no nulo \(\Psi=\Psi_\infty(x_\infty)\); una preparación estadística se
representa por un operador densidad positivo y de traza uno sobre la misma
realización. La identificación de fase global actúa en la representación:
la sección y su registro permanecen en el dominio genealógico.
\end{principle}

\subsection{Evolución, Hamiltoniano y ecuación de Schrödinger}
\label{v:v:rec:realizacion_evolucion:section:03}

La evolución genealógica preserva la composición de transportes. Su
realización satisface
\[
 U(t_3,t_2)U(t_2,t_1)=U(t_3,t_1),
 \qquad U(t,t)=I.
\]
Cuando preserva el producto interno y satisface
\(U(t)J_E=J_EU(t)\), es ortogonal en la carta real y unitario en la carta
compleja. Para un grupo unitario uniparamétrico fuertemente continuo, el
teorema de Stone proporciona un generador autoadjunto y \(J_E\)-lineal
\(H\)~\cite{v:Stone1932}:
\[
 U(t)=\exp\!\left(-J_E H t/\hbar_{\rm int}\right).
\]
Sobre el dominio \(D(H)\), la derivación temporal da
\begin{equation}
 \hbar_{\rm int}J_E\dot\Psi(t)=H\Psi(t),
 \label{v:v:rec:eq:schrodinger-genealogico-rev2}
\end{equation}
y la carta \(J_E\leftrightarrow i\) produce
\(i\hbar_{\rm int}\dot\Psi=H\Psi\).

El calendario \(C_{108}\) realiza esta correspondencia de manera finita:
el Hamiltoniano autoadjunto construido en su fuente tiene por
exponencial, durante \(t_0\), el desplazamiento unitario de una iteración. El
grupo continuo requiere adicionalmente continuidad fuerte; no se infiere
sólo de la periodicidad del calendario.
